\section{Discussion and Limitations}
\label{sec:discussion}

The main experiment substitutes reconstruction scores directly. It demonstrates sensitivity of the threshold computation to a writable calibration input, not the ability to generate IoT packets that produce the selected scores. The separate cross-device training-feature reservoir shows that model-scored benign feature rows outside the calibration split can supply effective donors; their availability or controllability by a real adversary is untested. The experimental trust boundary also assumes that the threshold code and final threshold cannot be overwritten directly while the input buffer can. A deployment would require explicit enforcement of that boundary.

With-replacement sampling from the same calibration tail reweights observed values and introduces detectable duplicates. The interpolated-score sensitivity preserves much of the measured displacement while reducing exact duplication, but it still does not establish traffic feasibility. The lowering attack uses random replacement positions; its effects should not be interpreted as the largest achievable under a rank-aware adversary.

Policy effects have different failure modes. \LOCAL confines the attack to its victim, whose detection loss can be large. \GLOBAL dilutes one local-threshold shift by $1/9$ yet changes all clients' operating points; its raw thresholds arise from heterogeneous client scales. \CLUSTER has assignment-dependent victim and non-victim effects, including opposite signs across clients. Fixed-assignment and $K$/initialization sensitivities limit any general claim that clustering reliably isolates contamination. The gate thresholds were chosen as consistency and materiality criteria, and sensitivity analyses show that gate decisions depend on those choices; percentile intervals at 20 seeds remain approximate.

The scope is one dataset, nine devices, a FedAvg autoencoder, $q{=}0.95$, and one compromised client at a time. Neither multi-client collusion nor calibration-input integrity mechanisms were evaluated. These limits constrain deployment claims while leaving the controlled score-level result intact.
